\documentclass[a4paper]{article}
\usepackage[pdftex]{graphicx}
\usepackage[utf8]{inputenc}
\usepackage{enumerate}
\usepackage{siunitx}
\sisetup{locale=DE}
\usepackage{href-ul}
\hypersetup{
	colorlinks=true,
	linkcolor=blue,
	urlcolor=blue}
\usepackage{icomma}
\usepackage{amssymb}
\usepackage{geometry}
\geometry{a4paper, top=15mm, left=15mm, right=15mm, bottom=15mm,
	headsep=10mm, footskip=12mm}

\begin{document}
	{\bf Accelerated movement}
	\begin{enumerate}[1.]
		
		%Task 1
		{\bf \item At the World Cup, the ball reached a speed of 108 km/h during a penalty kick.} If the goalkeeper had held the ball, the ball would have slowed down to zero speed within 0.04 s.
		
		\begin{enumerate}[a)]
			\item How large would the average, assumed constant deceleration (= negative acceleration) of the ball be when braking?
			\item Calculate the force that the goalkeeper would have had to apply for a ball with a mass of 420 g.
		\end{enumerate}
		
		%Exercise 2
		{\bf \item task}
		\begin{enumerate}[a)]
			\begin{minipage}{0.4\textwidth}
				\item Explain the movement of the t--v-diagram opposite in key words.
				\item From the diagram, explain that the distance traveled in phase 1 is twice as large as that in phase 2.
			\end{minipage}
			\hfill
			\begin{minipage}{0.45\textwidth}
				\includegraphics[width=5.5 cm]{vaute104.pdf}
			\end{minipage}
			\item Calculate the total distance the body travels in 12 s
			(Incorrect replacement result: 60 m)\\
			What is his average speed?
			\item The body whose movement the diagram describes has a mass of 750 g. In which of the 5 movement phases is the magnitude of the acting force greatest or smallest? How big are the forces in each case?
			\item The body now immediately (at t = 12 s) moves back the entire way. To do this, it starts with the acceleration $|\vec{a}|=\SI{4,0}{\meter\over {\second^2}}$ and in this way covers 32 m. He then covers the rest of the route at the same speed. How long does it take him to cover the entire distance?
			\item Qualitatively draw the t-v diagram of this return trip.
		\end{enumerate}
		
		{\bf \item A car is driving on the highway at a constant speed, the speedometer shows
			$\SI{125}{\kilo\meter\over \hour}$.} For a measurement, the time for a distance of $s_1=\SI{5.00}{\kilo\meter}$
		$t= \SI{2}{\minute}~\SI{30}{\second}$ stopped.
		\begin{enumerate}[a)]
			\item Calculate the actual speed $v$ in $ \SI{}{\meter\over \second} $ and
			$ \SI{}{\kilo\meter\over \hour} $
			\item What time does the car need at this speed for $s_2=\SI{5.40}{\kilo\meter}$
			long distance? Use the units $\SI{}{\minute}$ and $\SI{}{\second}$ in the answer.
			\item What distance would the car have traveled in $\SI{1.6}{\hour}$ if the speedometer reading was correct?
			\item Formulate an equivalent statement that it is a movement with constant speed.
			
		\end{enumerate}
		
		\begin{minipage}{0.65\textwidth}
			%Task 1
			{\bf \item A comparison by a motorsport magazine showed that of 7 super sports cars tested, the Porsche 911 Turbo S has the best acceleration values.} This car accelerates from 0 to $ \SI{200}{\kilo \meter \over \hour}$ in exactly
			$\SI{10,0}{\second}$.
		\end{minipage}
		\hfill
		\begin{minipage}{0.3\textwidth}
			\includegraphics[width=4 cm]{porsche129}
		\end{minipage}
		
		\begin{enumerate}[a)]
			\item Calculate the acceleration $a$ of the Porsche. It should be assumed that the acceleration $ a $ is constant.
			\item Determine the distance the car travels within this $\SI{10}{\second}$.
		\end{enumerate}
		
		
		{\bf During another test drive, the Porsche is brought from a standstill $(t_0=\SI{0}{\second})$ to $\SI{60}{\meter\over \second} $ within $ (t_1= \SI{15}{\second})$ accelerates.} From $ t_1=\SI{15}{\second}$ the sports car constantly maintains this speed until $ t_2=\SI{45}{\second}$, before accelerating within 20 seconds with constant acceleration to $\SI{90}{\meter\over \second} \quad (t_3=\SI{65}{\second})$.
		
		\begin{enumerate}[(i)]
			\item Draw a t-v diagram from time $t_0$ to time $t_3$ . Use a suitable scale.
			\item Use the diagram to determine the total distance traveled.
		\end{enumerate}
		
	\end{enumerate}
	
	\fbox{
		\begin{minipage}{0.5\textwidth}
			For the solution, please \href{https://www.okuyakl.de/phys/p10kinL410/le410.pdf}{click here} or scan the QR code.\\
			Additional worksheets are available at
			
			\href{http://www.okuyakl.com}{www.okuyakl.com}
		\end{minipage}
		\hfill
		\begin{minipage}{0.4\textwidth}
			\includegraphics[width=1.5 cm]{../../viecher/zwe03}
			\includegraphics[width=3 cm]{qre410}
			\includegraphics[width=2 cm]{../../viecher/afanticon1}
			
	\end{minipage}}
	
\end{document}